\documentclass[aspectratio=169,11pt]{beamer}

\usetheme{Madrid}
\usecolortheme{default}
\usefonttheme{professionalfonts}

\usepackage[T1]{fontenc}
\usepackage[utf8]{inputenc}
\usepackage[polish]{babel}
\usepackage{booktabs}
\usepackage{array}
\usepackage{hyperref}

\hyphenpenalty=10000
\exhyphenpenalty=10000
\righthyphenmin=62
\lefthyphenmin=62

\setbeamertemplate{navigation symbols}{}
\setbeamertemplate{footline}[frame number]

\title{Detekcja tekstu wygenerowanego przez AI}
\subtitle{Jak odróżnić esej człowieka od tekstu AI}
\author{Michał Dąbrowski \and Dominik Zabielski \and Michał Niechciałkowski}
\institute{Politechnika Białostocka, Wydział Informatyki\\
Prowadzący: dr hab.\ inż.\ Sławomir Zieliński}
\date{Białystok, 2026}

\begin{document}

\begin{frame}
  \titlepage
\end{frame}

\begin{frame}{Plan prezentacji}
  \begin{enumerate}
    \item O co chodzi w projekcie
    \item Z jakich danych korzystaliśmy
    \item Prosta metoda: cechy stylu tekstu
    \item DistilBERT i dlaczego wynik 99\% był mylący
    \item System One: lepszy model na esejach
    \item Aplikacja, test na nowych esejach i wnioski
  \end{enumerate}
\end{frame}

\begin{frame}{O co chodzi w projekcie}
  \begin{itemize}
    \item Chcemy rozpoznać, czy esej napisał człowiek, czy sztuczna inteligencja
    \item Metoda 1: proste cechy stylu tekstu (długość zdań, rytm, słownictwo)
    \item Metoda 2: modele językowe DistilBERT, a potem System One
    \item Sprawdzamy wyniki na tych samych danych, na obcych tekstach
          i na 300 ludzkich esejach z 2021~r.
    \item Patrzymy też, ile ludzkich tekstów mylnie uznajemy za AI
  \end{itemize}
\end{frame}

\begin{frame}{Zaczęliśmy od zbioru Kaggle}
  \begin{itemize}
    \item Angielskie eseje: 29\,145 tekstów, po czyszczeniu 27\,296
    \item Usunęliśmy powtórzenia; ok.\ 60\% człowiek, 40\% AI
  \end{itemize}

  \vspace{0.6em}
  \begin{center}
  \begin{tabular}{lr}
    \toprule
    podział & teksty \\
    \midrule
    nauka modelu & 21\,838 \\
    sprawdzanie w trakcie & 2\,729 \\
    test końcowy & 2\,729 \\
    \bottomrule
  \end{tabular}
  \end{center}

  \vspace{0.8em}
  \alert{To był tylko początek.} Wynik na samym Kaggle okazał się mylący.
\end{frame}

\begin{frame}{Metoda statystyczna}
  \begin{itemize}
    \item 13 prostych cech z tekstu: długość zdań, skokowość rytmu,
          bogactwo słów, przecinki
    \item Regresja logistyczna: ok.\ 91\% trafień, 254 błędy na 2\,728 tekstów
    \item Za człowiekiem: dłuższe i bardziej nierówne zdania
    \item Za AI: dłuższe słowa i więcej przecinków
    \item Plus: decyzje da się wyjaśnić
    \item Minus: na krótkich tekstach działa gorzej, a promptem da się ją oszukać
  \end{itemize}
\end{frame}

\begin{frame}{Wynik 99\% był mylący}
  DistilBERT uczony tylko na Kaggle wygląda świetnie,
  dopóki nie wyjdziemy poza ten zbiór.

  \vspace{1em}
  \begin{center}
  \begin{tabular}{lcc}
    \toprule
    & Kaggle & obce teksty Ghostbuster \\
    \midrule
    DistilBERT (tylko Kaggle) & \textbf{99,7\%} & \alert{51,6\%} \\
    \bottomrule
  \end{tabular}
  \end{center}

  \vspace{1em}
  Na nowych danych model prawie zawsze mówi ,,to AI''.

  \vspace{0.5em}
  \textbf{Wniosek:} jeden zbiór z Kaggle nie pokazuje,
  jak detektor działa naprawdę.
\end{frame}

\begin{frame}{Trening na wielu źródłach}
  \begin{itemize}
    \item Do nauki dodaliśmy teksty z kilku zbiorów, nie tylko z Kaggle
    \item Osobno sprawdzamy Ghostbustera, którego model wcześniej nie widział
  \end{itemize}

  \vspace{0.8em}
  \begin{center}
  \begin{tabular}{lc}
    \toprule
    model na obcych tekstach & trafność \\
    \midrule
    stary model (tylko Kaggle) & 52\% \\
    DistilBERT na wielu źródłach & \textbf{84\%} \\
    \bottomrule
  \end{tabular}
  \end{center}

  \vspace{0.8em}
  Duży skok, ale ludzkie teksty nadal bywają mylone z AI.
  Dlatego zbudowaliśmy System One.
\end{frame}

\begin{frame}{System One: jak działa}
  Model dostaje esej i porównuje dwie możliwości:

  \vspace{0.6em}
  \begin{columns}[T]
    \begin{column}{0.48\textwidth}
      \begin{block}{opcja: człowiek}
        nierówny rytm, szczegóły, głos studenta
      \end{block}
    \end{column}
    \begin{column}{0.48\textwidth}
      \begin{block}{opcja: AI}
        szablon, wygładzony styl, ogólnikowa teza
      \end{block}
    \end{column}
  \end{columns}

  \vspace{0.6em}
  \begin{itemize}
    \item Z wyników obu opcji powstaje prawdopodobieństwo ,,to AI''
    \item Nauka na wielu zbiorach esejów (ok.\ 52 tys.\ tekstów)
    \item Trening ok.\ 37 minut; gotowy model: \texttt{models/system-one-final/}
  \end{itemize}
\end{frame}

\begin{frame}{System One: wyniki}
  Na obcych tekstach jest wyraźnie lepszy
  i rzadziej oskarża człowieka o AI.

  \vspace{0.8em}
  \begin{center}
  \begin{tabular}{lccc}
    \toprule
    model & obce teksty & ludzkie OK & AI wykryte \\
    \midrule
    DistilBERT (wiele źródeł) & 0,838 & 0,705 & 0,970 \\
    System One v3 (finalny) & \textbf{0,946} & \textbf{0,900} & 0,993 \\
    \bottomrule
  \end{tabular}
  \end{center}

  \vspace{0.8em}
  Mylne alarmy na Ghostbusterze spadły z 279 do 30.
\end{frame}

\begin{frame}{Aplikacja i test na nieznanych esejach}
  \begin{itemize}
    \item Model widzi tylko początek tekstu, zwykle wstęp
    \item Aplikacja dzieli esej na fragmenty, ocenia każdy i bierze średnią
    \item Test: 300 ludzkich esejów z 2021~r., sprzed ChatGPT, nigdy w treningu
  \end{itemize}

  \vspace{0.6em}
  \begin{center}
  \begin{tabular}{lccc}
    \toprule
    detektor & człowiek & niepewne & mylnie AI \\
    \midrule
    stara aplikacja & -- & -- & 93,3\% \\
    DistilBERT z Kaggle & -- & -- & 96,3\% \\
    metoda statystyczna & -- & -- & 18,3\% \\
    System One, fragmenty & \textbf{80,3\%} & 12,7\% & \textbf{7,0\%} \\
    \bottomrule
  \end{tabular}
  \end{center}
\end{frame}

\begin{frame}{Czego jeszcze nie rozwiązaliśmy}
  \begin{itemize}
    \item Nowe modele AI i przerobione teksty nadal bywają trudne
    \item Początek i koniec ludzkiego eseju czasem wygląda jak AI
    \item Na esejach z 2021~r.\ nadal 7\% mylnych alarmów
          i ok.\ 13\% odpowiedzi ,,niepewne''
    \item Bardzo gładki tekst nadal skłania model do odpowiedzi ,,AI''
    \item Statystykę łatwo oszukać poleceniem
          ,,pisz zdaniami różnej długości''
  \end{itemize}
\end{frame}

\begin{frame}{Podsumowanie}
  \begin{itemize}
    \item Sam Kaggle zawyża wynik: 99,7\% potrafi spaść do ok.\ 52\%
    \item Statystyka: ok.\ 91\% na Kaggle, łatwa do wyjaśnienia,
          słabsza poza tym zbiorem
    \item System One: ok.\ 95\% na obcych tekstach
    \item W aplikacji mylne ,,AI'' na esejach z 2021~r.\ spadło z 93\% do 7\%
    \item Na końcu używamy System One, a nie pierwszego modelu z Kaggle
  \end{itemize}
\end{frame}

\begin{frame}
  \centering
  \vspace{1.5cm}
  {\LARGE Dziękujemy za uwagę}

  \vspace{1cm}
  {\Large Pytania?}

  \vspace{1.2cm}
  Michał Dąbrowski · Dominik Zabielski · Michał Niechciałkowski
\end{frame}

\end{document}
